\documentclass[12pt]{article}
\usepackage[utf8]{inputenc}
\usepackage[T1]{fontenc}
\usepackage[margin=1in]{geometry}
\usepackage{setspace}
\usepackage{xcolor}
\usepackage[colorlinks=true, urlcolor=blue!60!black, linkcolor=blue!60!black]{hyperref}
\setlength{\parindent}{2em}
\setlength{\parskip}{1.25em}
\renewcommand{\baselinestretch}{1.0}
\begin{document}
\begin{center}
{\LARGE \textit{Texas v. Johnson}, 491 U.S. 397 (1989)\par}
\vspace{0.8em}
{\large Case Brief\par}
\end{center}
\vspace{1em}

\section{Memory Jogger}
Burning an American flag as political protest is expressive conduct protected by the First Amendment, and a State may not criminalize that act simply because society finds the message offensive or wishes to preserve the flag as a symbol of national unity (\href{https://supreme.justia.com/cases/federal/us/491/397/}{Justia}; \href{https://www.courtlistener.com/opinion/112304/texas-v-johnson/}{CourtListener}).

\section{Detailed Case Facts}
During the 1984 Republican National Convention in Dallas, Texas, Gregory Lee Johnson participated in a political demonstration dubbed the ``Republican War Chest Tour,'' organized to protest the policies of the Reagan administration and certain Dallas-based corporations (\href{https://supreme.justia.com/cases/federal/us/491/397/}{491 U.S. at 399}). Roughly 100 demonstrators marched through Dallas streets, chanted political slogans, distributed literature, and staged ``die-ins'' at corporate offices to dramatize the threat of nuclear war (\href{https://supreme.justia.com/cases/federal/us/491/397/}{491 U.S. at 399--400}). Although some protesters spray-painted walls and overturned potted plants along the route, Johnson himself did not participate in those acts (\href{https://supreme.justia.com/cases/federal/us/491/397/}{491 U.S. at 399}). Another demonstrator removed an American flag from a flagpole outside one of the targeted buildings and handed it to Johnson (\href{https://supreme.justia.com/cases/federal/us/491/397/}{491 U.S. at 399}).

When the march reached Dallas City Hall, Johnson unfurled the American flag, doused it with kerosene, and set it ablaze while fellow protesters chanted, ``America, the red, white, and blue, we spit on you'' (\href{https://supreme.justia.com/cases/federal/us/491/397/}{491 U.S. at 399}). No one was physically injured or threatened with physical injury during the incident, though several onlookers later testified that they were seriously offended (\href{https://supreme.justia.com/cases/federal/us/491/397/}{491 U.S. at 399}). After the crowd dispersed, a witness gathered the burned remains of the flag and buried them in his backyard (\href{https://supreme.justia.com/cases/federal/us/491/397/}{491 U.S. at 399}).

Of the approximately 100 demonstrators, Johnson alone was criminally charged, and the sole offense prosecuted was desecration of a venerated object under Texas Penal Code Ann. \S~42.09(a)(3) (1989) (\href{https://supreme.justia.com/cases/federal/us/491/397/}{491 U.S. at 400 \& n.1}). That statute made it a Class~A misdemeanor to intentionally or knowingly desecrate a state or national flag, defining ``desecrate'' as to ``deface, damage, or otherwise physically mistreat in a way that the actor knows will seriously offend one or more persons likely to observe or discover his action'' (\href{https://supreme.justia.com/cases/federal/us/491/397/}{491 U.S. at 400 n.1}).

\section{Procedural History}
After a jury trial in Texas state court, Johnson was convicted of violating Texas Penal Code Ann. \S~42.09(a)(3), sentenced to one year in prison, and fined \$2,000 (\href{https://supreme.justia.com/cases/federal/us/491/397/}{491 U.S. at 400}). The Court of Appeals for the Fifth District of Texas at Dallas affirmed his conviction, 706 S.W.2d 120 (Tex. App. 1986) (\href{https://supreme.justia.com/cases/federal/us/491/397/}{491 U.S. at 400}).

The Texas Court of Criminal Appeals reversed, 755 S.W.2d 92 (Tex. Crim. App. 1988), holding that \S~42.09 was unconstitutional under the First Amendment as applied to Johnson's political protest (\href{https://supreme.justia.com/cases/federal/us/491/397/}{491 U.S. at 400--402}). That court concluded that Johnson's flag burning was protected symbolic speech, that the State could not mandate national unity by prescribing approved messages for the flag under \textit{West Virginia State Board of Education v. Barnette}, 319 U.S. 624 (1943), and that no actual or threatened breach of the peace occurred (noting that Texas already had a separate breach-of-the-peace statute, Tex. Penal Code Ann. \S~42.01) (\href{https://supreme.justia.com/cases/federal/us/491/397/}{491 U.S. at 400--401}). The U.S. Supreme Court granted certiorari, 488 U.S. 907 (1988), heard oral argument on March 21, 1989, and affirmed on June 21, 1989 (\href{https://supreme.justia.com/cases/federal/us/491/397/}{491 U.S. at 397, 402}; \href{https://www.oyez.org/cases/1988/88-155}{Oyez}).

\section{Judicial Votes}
\textbf{Vote Split:} 5--4, affirming the Texas Court of Criminal Appeals (\href{https://supreme.justia.com/cases/federal/us/491/397/}{491 U.S. at 398}; \href{https://constitutioncenter.org/the-constitution/supreme-court-case-library/texas-v-johnson}{National Constitution Center}).

\begin{itemize}
    \item \textbf{Majority (5):} Justice William J. Brennan, Jr. delivered the opinion of the Court, joined by Justices Thurgood Marshall, Harry A. Blackmun, Antonin Scalia, and Anthony M. Kennedy.
    \item \textbf{Concurrence:} Justice Anthony M. Kennedy filed a concurring opinion (491 U.S. at 420), joining Justice Brennan's opinion in full without reservation.
    \item \textbf{Dissents (4):} Chief Justice William H. Rehnquist filed a dissenting opinion, joined by Justices Byron R. White and Sandra Day O'Connor (491 U.S. at 421). Justice John Paul Stevens filed a separate dissenting opinion (491 U.S. at 436).
\end{itemize}

\section{Holding}
\textbf{Judgment:} Affirmed. Johnson's criminal conviction for flag desecration under Texas Penal Code Ann. \S~42.09(a)(3) was reversed as inconsistent with the First Amendment (\href{https://supreme.justia.com/cases/federal/us/491/397/}{491 U.S. at 399, 420}).

\textbf{Rule Announced:} First Amendment protection extends beyond spoken and written words to expressive conduct where an intent to convey a particularized message is present and the likelihood is great that the message will be understood by viewers (\textit{Spence v. Washington}, 418 U.S. 405 (1974)) (\href{https://supreme.justia.com/cases/federal/us/491/397/}{491 U.S. at 404}). When the State regulates expressive conduct for reasons related to the suppression of expression---such as protecting a symbol's official meaning or shielding onlookers from serious offense---the lenient intermediate scrutiny of \textit{United States v. O'Brien}, 391 U.S. 367 (1968), does not apply; instead, the content-based restriction triggers ``the most exacting scrutiny'' (\textit{Boos v. Barry}, 485 U.S. 312 (1988)), which cannot be satisfied by a desire to prohibit offensive ideas or enforce a single orthodox meaning for the American flag (\href{https://supreme.justia.com/cases/federal/us/491/397/}{491 U.S. at 407, 412, 414--417}).

\textbf{Holding vs. Dicta:} The Court resolved the case strictly on as-applied grounds, holding that \S~42.09 violated the First Amendment as applied to Johnson's political protest, and expressly declined to adjudicate Johnson's facial vagueness and overbreadth challenges (\href{https://supreme.justia.com/cases/federal/us/491/397/}{491 U.S. at 402, 403 n.3}). The majority also noted in footnote~8 that its decision did not immunize theft of a flag, trespass, disorderly conduct, or arson, because Johnson was prosecuted solely for flag desecration (\href{https://supreme.justia.com/cases/federal/us/491/397/}{491 U.S. at 412 n.8}). Similarly, the majority's remarks that tolerating flag burning strengthens the flag's cherished place in American life and that saluting a burning flag or burying its remains is the most fitting response constitute persuasive rhetorical dicta reinforcing the Court's First Amendment principles (\href{https://supreme.justia.com/cases/federal/us/491/397/}{491 U.S. at 419--420}).

\section{Analysis of Opinions}

\subsection{Majority Opinion (Justice Brennan)}
Justice Brennan structured the opinion in three sequential steps (\href{https://supreme.justia.com/cases/federal/us/491/397/}{491 U.S. at 403--404}). First, applying the \textit{Spence} test, the Court held that Johnson's flag burning was ``sufficiently imbued with elements of communication'' to implicate the First Amendment (\href{https://supreme.justia.com/cases/federal/us/491/397/}{491 U.S. at 404--406}). Texas conceded the point at oral argument, and the context---the culmination of a political march timed to coincide with President Reagan's renomination---made the political message ``both intentional and overwhelmingly apparent'' (\href{https://supreme.justia.com/cases/federal/us/491/397/}{491 U.S. at 405--406}).

Second, Justice Brennan examined whether either of Texas's two asserted interests was ``unrelated to the suppression of free expression'' so as to trigger the more lenient \textit{O'Brien} standard (\href{https://supreme.justia.com/cases/federal/us/491/397/}{491 U.S. at 406--407}). As to preventing breaches of the peace, the Court held that interest was simply not implicated on the record: no disturbance actually occurred or threatened to occur, and the State cannot presume that provocative or offensive speech will incite violence without meeting the imminent-lawless-action test of \textit{Brandenburg v. Ohio}, 395 U.S. 444 (1969), or the direct personal-insult requirement of ``fighting words'' under \textit{Chaplinsky v. New Hampshire}, 315 U.S. 568 (1942) (\href{https://supreme.justia.com/cases/federal/us/491/397/}{491 U.S. at 407--410}). As to preserving the flag as a symbol of nationhood and national unity, that interest ``blossom[s] only when a person's treatment of the flag communicates some message'' and is therefore directly related to the suppression of expression, taking the case outside \textit{O'Brien} (\href{https://supreme.justia.com/cases/federal/us/491/397/}{491 U.S. at 410}).

Third, applying ``the most exacting scrutiny'' because \S~42.09 turned on the communicative impact of the act---punishing flag burning only when it causes ``serious offense'' while permitting respectful burning to dispose of a worn flag under 36 U.S.C. \S~176(k)---the Court held the conviction unconstitutional (\href{https://supreme.justia.com/cases/federal/us/491/397/}{491 U.S. at 411--412}). Invoking \textit{Barnette}, \textit{Street v. New York}, 394 U.S. 576 (1969), and \textit{Schacht v. United States}, 398 U.S. 58 (1970), Justice Brennan declared that the government may not prohibit expression simply because society finds it offensive, nor prescribe a single orthodox meaning for designated national symbols (\href{https://supreme.justia.com/cases/federal/us/491/397/}{491 U.S. at 414--418}). Creating a flag-only exception would have ``no discernible or defensible boundaries''---inviting bans on burning state flags, the Presidential seal, or copies of the Constitution---and would dilute the very freedom the flag represents (\href{https://supreme.justia.com/cases/federal/us/491/397/}{491 U.S. at 417--420}).

\subsection{Concurring Opinion (Justice Kennedy)}
Justice Kennedy joined Justice Brennan's opinion ``without reservation'' but wrote a brief, candid concurrence acknowledging the ``personal toll'' of the judgment (\href{https://supreme.justia.com/cases/federal/us/491/397/}{491 U.S. at 420}). Emphasizing judicial duty over personal preference, he observed that ``sometimes we must make decisions we do not like'' because the Constitution compels the result (\href{https://supreme.justia.com/cases/federal/us/491/397/}{491 U.S. at 420--421}). While honoring the flag's unique place for those who carried it in battle, Justice Kennedy concluded that ``[i]t is poignant but fundamental that the flag protects those who hold it in contempt'' (\href{https://supreme.justia.com/cases/federal/us/491/397/}{491 U.S. at 421}).

\subsection{Dissenting Opinion (Chief Justice Rehnquist)}
Chief Justice Rehnquist, joined by Justices White and O'Connor, invoked Justice Holmes's aphorism that ``a page of history is worth a volume of logic,'' recounting over 200 years of American history, poetry (including Emerson's \textit{Concord Hymn}, Francis Scott Key's \textit{The Star-Spangled Banner}, and Whittier's \textit{Barbara Frietchie}), wartime sacrifice at Iwo Jima and Inchon, and federal and state legislation---including flag-desecration laws in 48 of the 50 States---to demonstrate the flag's unique status as the visible symbol of the Nation (\href{https://supreme.justia.com/cases/federal/us/491/397/}{491 U.S. at 421--429}).

Doctrinally, Chief Justice Rehnquist analogized flag burning to ``fighting words'' under \textit{Chaplinsky}, characterizing it as ``the equivalent of an inarticulate grunt or roar'' of slight social value as a step to truth that tends to incite a breach of the peace (\href{https://supreme.justia.com/cases/federal/us/491/397/}{491 U.S. at 430--432}). He argued that \S~42.09 did not punish Johnson's anti-Reagan ideas---leaving him free to march, chant, and burn effigies---but merely prohibited one inflammatory mode of expression (\href{https://supreme.justia.com/cases/federal/us/491/397/}{491 U.S. at 431--432}). Relying on \textit{Halter v. Nebraska}, 205 U.S. 34 (1907), and \textit{San Francisco Arts \& Athletics, Inc. v. United States Olympic Committee}, 483 U.S. 522 (1987), he concluded that the Court acted as a ``platonic guardian'' improperly overriding democratic majorities (\href{https://supreme.justia.com/cases/federal/us/491/397/}{491 U.S. at 429--430, 435}).

\subsection{Dissenting Opinion (Justice Stevens)}
Justice Stevens dissented separately, framing the flag as an intangible ``important national asset'' symbolizing not merely national unity, but the ideals of freedom, equal opportunity, and tolerance (\href{https://supreme.justia.com/cases/federal/us/491/397/}{491 U.S. at 436--439}). Analogizing public flag burning to spray-painting graffiti or posting bulletin boards on the Washington Monument or Lincoln Memorial, he argued that sanctioning flag desecration tarnishes the symbol's value at a cost he would not pay, while imposing only a ``trivial burden'' on speakers who retain ample alternative avenues to criticize the government (\href{https://supreme.justia.com/cases/federal/us/491/397/}{491 U.S. at 437--439}). In his view, \S~42.09 was content-neutral because it penalized the prohibited physical method of desecration regardless of the speaker's underlying political viewpoint (\href{https://supreme.justia.com/cases/federal/us/491/397/}{491 U.S. at 438--439 \& n.*}).

\section{Examples: Future Applications}
\begin{enumerate}
    \item \textbf{Same-Side (Protected Expression):} Congress enacts a facially neutral federal statute---the Flag Protection Act of 1989---criminalizing anyone who ``knowingly mutilates, defaces, physically defiles, burns, maintains on the floor or ground, or tramples upon any flag of the United States,'' omitting any explicit requirement of ``serious offense.'' As the Supreme Court confirmed in \textit{United States v. Eichman}, 496 U.S. 310 (1990), the statute still violates the First Amendment under \textit{Johnson} because the government's asserted interest in protecting the flag's symbolic meaning remains directly related to the suppression of expression (\href{https://supreme.justia.com/cases/federal/us/491/397/}{Justia}; \href{https://illinoislawreview.org/wp-content/uploads/2022/10/Flanders.pdf}{Illinois Law Review Online}).

    \item \textbf{Same-Side (Protected Expression):} A state legislature passes a law prohibiting the public burning or defacement of copies of the state flag or the state constitution ``in a manner likely to cause public outrage.'' Because the law targets the communicative impact of political dissent directed at a designated government symbol, \textit{Johnson}'s anti-orthodoxy principle triggers strict scrutiny and invalidates the law (\href{https://supreme.justia.com/cases/federal/us/491/397/}{491 U.S. at 417}).

    \item \textbf{Opposite-Side (Permissible Prosecution):} A political protester steals an American flag from a private homeowner's porch, trespasses onto city property, and ignites it in violation of a generally applicable municipal fire-safety ordinance banning open flames in a dry public park. Under \textit{Johnson} footnote~8 and \textit{O'Brien}, the State may prosecute the protester for theft, trespass, or violating the content-neutral open-flame ordinance because those laws serve substantial nonspeech interests (property rights and fire prevention) unrelated to the suppression of expression (\href{https://supreme.justia.com/cases/federal/us/491/397/}{491 U.S. at 407, 412 n.8}).

    \item \textbf{Opposite-Side (Permissible Prosecution):} A demonstrator spray-paints anti-war slogans onto the marble facade of the Lincoln Memorial. Even though the message is core political speech, the government may enforce a content-neutral ban on defacing public monuments because protecting physical government property is unrelated to suppressing any particular viewpoint (\href{https://chicagounbound.uchicago.edu/cgi/viewcontent.cgi?article=1010&context=occasional_papers}{Geoffrey R. Stone, Chicago Unbound}).

    \item \textbf{Fence-Sitter (Close Call):} A protester burns a privately owned American flag while standing face-to-face inches from a grieving military family at a cemetery entrance, shouting personal threats and urging a hostile surrounding crowd to attack bystanders. Although burning the flag itself is protected symbolic speech under \textit{Johnson}, the surrounding conduct and direct personal confrontation push the incident to the edge of unprotected ``fighting words'' (\textit{Chaplinsky}) or incitement to ``imminent lawless action'' (\textit{Brandenburg}), making constitutionality turn heavily on the specific factual record (\href{https://supreme.justia.com/cases/federal/us/491/397/}{491 U.S. at 409--410}).
\end{enumerate}

\section{Critique}
\textbf{Originalist and Textualist Perspectives:} \textit{Texas v. Johnson} presents a famous puzzle for originalists and textualists. On the one hand, Justice Scalia joined Justice Brennan's majority opinion in full---later pointing to \textit{Johnson} as proof that principled adherence to constitutional rules against viewpoint-based censorship overrides personal political distaste for ``flag-burning'' protesters (\href{https://scholar.valpo.edu/cgi/viewcontent.cgi?article=1450&context=law_fac_pubs}{Valparaiso Law Review}; \href{https://firstamendment.mtsu.edu/article/antonin-scalia/}{First Amendment Encyclopedia}). On the other hand, strict textualists and history-and-tradition originalists---including Justice Hugo Black in \textit{Street v. New York} and contemporary scholars examining \textit{Johnson} after \textit{New York State Rifle \& Pistol Association v. Bruen}---have questioned whether nonverbal conduct like burning a physical object falls within the original public meaning of ``speech'' and noted that Brennan's opinion relied almost exclusively on modern doctrinal precedents rather than Founding-era text or history (\href{https://illinoislawreview.org/wp-content/uploads/2022/10/Flanders.pdf}{Chad Flanders, Illinois Law Review Online}; \href{https://scholar.law.colorado.edu/cgi/viewcontent.cgi?article=1739&context=faculty-articles}{Paul Campos, California Law Review}).

\textbf{Progressive and Doctrinal Perspectives:} Mainstream First Amendment scholars---including Geoffrey R. Stone and John Hart Ely---widely praise \textit{Johnson} as a triumph of principled doctrinal coherence (\href{https://chicagounbound.uchicago.edu/cgi/viewcontent.cgi?article=1010&context=occasional_papers}{Geoffrey R. Stone, Chicago Unbound}; \href{https://lawpublications.barry.edu/cgi/viewcontent.cgi?article=1063&context=barrylrev}{Barry Law Review}). Building on Ely's influential analysis of \textit{O'Brien}, Stone argued that \textit{Johnson} correctly enforced the foundational line between content-neutral regulations of conduct and content-based restrictions triggered by audience offense, thereby preventing majorities from carving out ad hoc ideological exceptions to the marketplace of ideas (\href{https://chicagounbound.uchicago.edu/cgi/viewcontent.cgi?article=1010&context=occasional_papers}{Geoffrey R. Stone, Chicago Unbound}). Some critics, however, contend that the majority's insistence on strict content neutrality undervalues civic cohesion and treats a uniquely unifying national emblem identically to ordinary commercial or partisan symbols (\href{https://supreme.justia.com/cases/federal/us/491/397/}{491 U.S. at 429, 437}).

\textbf{Competing Constitutional Values:} \textit{Johnson} crystallizes the enduring tension between two foundational democratic ideals: (1)~radical individual liberty of dissent, anti-orthodoxy, and viewpoint neutrality (forbidding the State to punish deeply offensive political protest), and (2)~communitarian self-governance and civic solidarity (allowing democratic majorities to safeguard a shared national symbol that binds a diverse republic together) (\href{https://supreme.justia.com/cases/federal/us/491/397/}{491 U.S. at 414--420, 435}).

\section{Key Quotations}
\begin{quote}
``If there is a bedrock principle underlying the First Amendment, it is that the government may not prohibit the expression of an idea simply because society finds the idea itself offensive or disagreeable.''
\end{quote}
--- Justice William J. Brennan, Jr., Opinion of the Court, 491 U.S. at 414 (\href{https://supreme.justia.com/cases/federal/us/491/397/}{Justia}; \href{https://www.courtlistener.com/opinion/112304/texas-v-johnson/}{CourtListener}).

\begin{quote}
``We do not consecrate the flag by punishing its desecration, for in doing so we dilute the freedom that this cherished emblem represents.''
\end{quote}
--- Justice William J. Brennan, Jr., Opinion of the Court, 491 U.S. at 420 (\href{https://supreme.justia.com/cases/federal/us/491/397/}{Justia}).

\begin{quote}
``The hard fact is that sometimes we must make decisions we do not like. We make them because they are right, right in the sense that the law and the Constitution, as we see them, compel the result.''
\end{quote}
--- Justice Anthony M. Kennedy, Concurring Opinion, 491 U.S. at 420--421 (\href{https://supreme.justia.com/cases/federal/us/491/397/}{Justia}).

\begin{quote}
``In holding this Texas statute unconstitutional, the Court ignores Justice Holmes' familiar aphorism that `a page of history is worth a volume of logic.' For more than 200 years, the American flag has occupied a unique position as the symbol of our Nation, a uniqueness that justifies a governmental prohibition against flag burning in the way respondent Johnson did here.''
\end{quote}
--- Chief Justice William H. Rehnquist, Dissenting Opinion, 491 U.S. at 421--422 (\href{https://supreme.justia.com/cases/federal/us/491/397/}{Justia}).

\begin{quote}
``The creation of a federal right to post bulletin boards and graffiti on the Washington Monument might enlarge the market for free expression, but at a cost I would not pay.''
\end{quote}
--- Justice John Paul Stevens, Dissenting Opinion, 491 U.S. at 437 (\href{https://supreme.justia.com/cases/federal/us/491/397/}{Justia}).

\end{document}
